\chapter{Fundamentação Teórica e Trabalhos Correlatos}
\label{cap:fundamentacao}

O desenvolvimento de uma plataforma computacional para simulação física e experimentação interativa de sistemas de controle requer uma base teórica multidisciplinar. Este capítulo estabelece os fundamentos matemáticos e físicos da teoria de controle clássico por realimentação, detalha as ações do algoritmo PID, as não-linearidades de atuadores e os métodos empíricos de sintonia. Em seguida, revisam-se os princípios pedagógicos dos laboratórios virtuais no ensino de engenharia, os pilares computacionais adotados na implementação do sistema e um quadro comparativo com trabalhos correlatos.

\section{Sistemas de Controle em Malha Fechada}
\label{sec:controle_malha_fechada}

Em sistemas dinâmicos, a estratégia de controle em malha fechada (\textit{feedback control}) baseia-se na medição contínua da variável de processo $y(t)$ e em sua comparação com o valor de referência desejado, denominado \textit{setpoint} $r(t)$ \cite{ogata2010}. O desvio instantâneo entre a referência e a saída mensurada define o sinal de erro $e(t)$:

\begin{equation}
    e(t) = r(t) - y(t)
    \label{eq:sinal_erro}
\end{equation}

A função precípua do controlador $C(s)$ é processar dinamicamente o sinal de erro $e(t)$ para sintetizar um sinal de atuação ou esforço de controle $u(t)$, aplicado à planta $G(s)$ de modo a compelir a variável de processo a rastrear a referência e rejeitar perturbações externas $d(t)$, conforme ilustrado no diagrama de blocos canônico da Figura \ref{fig:diagrama_malha_fechada}.

\begin{figure}[htbp]
    \centering
    \caption{Diagrama de blocos canônico de um sistema de controle em malha fechada com realimentação unitária e perturbação aditiva.}
    \label{fig:diagrama_malha_fechada}
    \fbox{\parbox{0.85\textwidth}{\centering
        \vspace{0.4cm}
        $[r(t)] \xrightarrow{+} \bigoplus \xrightarrow{e(t)} \boxed{C(s)} \xrightarrow{u(t)} \bigoplus \xrightarrow{} \boxed{G(s)} \xrightarrow{y(t)} \bullet \xrightarrow{} [y(t)]$ \\
        \vspace{0.1cm}
        \hspace{6.2cm} $\uparrow +$ \hspace{2.8cm} $\vert$ \\
        \hspace{6.2cm} $[d(t)]$ \hspace{2.5cm} $\downarrow$ \\
        \hspace{2.2cm} $\xleftarrow[\hspace{7.5cm}]{-} \bullet \longleftarrow$
        \vspace{0.4cm}
    }}
    \vspace{0.1cm}
    \fonte{Elaborada pelo autor (2026).}
\end{figure}

\subsection{Ações do Controlador Proporcional-Integral-Derivativo (PID)}
\label{subsec:acoes_pid}

O algoritmo PID sintetiza o sinal de controle $u(t)$ a partir de uma combinação linear de três termos funcionais que atuam sobre o erro passado, presente e futuro estimado \cite{astrom2006, nise2020}. Na formulação clássica contínua em paralelo, o sinal de controle é definido por:

\begin{equation}
    u(t) = u_P(t) + u_I(t) + u_D(t)
    \label{eq:pid_soma}
\end{equation}

\noindent onde as parcelas individuais correspondem a:

\begin{align}
    u_P(t) &= K_p \, e(t) \label{eq:termo_p} \\
    u_I(t) &= K_i \int_{0}^{t} e(\tau) \, d\tau \label{eq:termo_i} \\
    u_D(t) &= K_d \, \frac{de(t)}{dt} \label{eq:termo_d}
\end{align}

Nesta parametrização, $K_p$ denota o ganho proporcional, $K_i$ o ganho integral e $K_d$ o ganho derivativo. Alternativamente, expressa-se a equação na forma padrão (ideal), parametrizada pelo tempo integral $T_i$ e pelo tempo derivativo $T_d$:

\begin{equation}
    u(t) = K_p \left( e(t) + \frac{1}{T_i} \int_{0}^{t} e(\tau) \, d\tau + T_d \, \frac{de(t)}{dt} \right)
    \label{eq:pid_padrao}
\end{equation}

\noindent sendo as relações diretas dadas por $K_i = \frac{K_p}{T_i}$ e $K_d = K_p T_d$.

Aplicando-se a Transformada de Laplace sob condições iniciais nulas, obtém-se a função de transferência do controlador no domínio da frequência complexa $s$:

\begin{equation}
    C(s) = \frac{U(s)}{E(s)} = K_p + \frac{K_i}{s} + K_d \, s = K_p \left( 1 + \frac{1}{T_i s} + T_d s \right)
    \label{eq:pid_laplace}
\end{equation}

Cada componente desempenha um papel físico e dinâmico bem definido na estabilização do sistema:

\begin{enumerate}
    \item \textbf{Ação Proporcional ($u_P$):} Gera uma força diretamente proporcional à magnitude do erro instantâneo. Aumenta a velocidade de resposta da malha fechada e reduz o erro de regime, contudo, ganhos $K_p$ excessivos induzem oscilações acentuadas e podem desestabilizar o sistema. Além disso, atuando isoladamente em sistemas de ordem zero ou tipo zero, a ação proporcional é incapaz de eliminar completamente o erro de regime permanente ($e_{ss} \neq 0$).
    \item \textbf{Ação Integral ($u_I$):} Acumula o histórico do erro ao longo do tempo. Sua inclusão adiciona um polo na origem ($s = 0$) à malha aberta, elevando o tipo do sistema e garantindo erro nulo em regime permanente ($e_{ss} = 0$) para entradas do tipo degrau e rejeição a perturbações constantes. Todavia, a ação integral introduz defasagem negativa no sistema, o que tende a degradar a margem de estabilidade e provocar sobressinais expressivos.
    \item \textbf{Ação Derivativa ($u_D$):} Responde à taxa de variação temporal do erro, atuando como um mecanismo preditivo ou amortecedor dinâmico. Ao antecipar a tendência de crescimento do erro, o termo derivativo adiciona amortecimento relativo ao sistema, mitigando oscilações e reduzindo o tempo de acomodação. Em contrapartida, a diferenciação pura é inerentemente sensível a ruídos de alta frequência presentes nos sinais de medição.
\end{enumerate}

\subsubsection{Discretização Numérica do Controlador}
\label{subsubsec:discretizacao_pid}

Em aplicações computacionais e sistemas embarcados, a lei de controle contínua deve ser discretizada para execução a intervalos regulares de amostragem $\Delta t$. Empregando-se a aproximação de diferenças finitas regressivas (\textit{backward differences}) e integração numérica retangular à esquerda (Euler explícito), a formulação discreta para o instante de amostragem $k$ resulta em:

\begin{align}
    u_P[k] &= K_p \, e[k] \label{eq:p_discreto} \\
    u_I[k] &= u_I[k-1] + K_i \, e[k] \, \Delta t \label{eq:i_discreto} \\
    u_D[k] &= K_d \, \frac{e[k] - e[k-1]}{\Delta t} \label{eq:d_discreto} \\
    u[k] &= u_P[k] + u_I[k] + u_D[k] \label{eq:pid_discreto}
\end{align}

Esta representação posicional direta foi implementada no motor de execução do \textit{PID Simulator}, garantindo que a execução do código escrito pelo discente no navegador preserve correspondência exata com o equacionamento discreto.

\subsection{Métricas de Desempenho Transitório e de Regime Permanente}
\label{subsec:metricas_desempenho}

A avaliação quantitativa da eficácia de um controlador baseia-se na resposta ao degrau unitário do sistema em malha fechada \cite{dorf2017}. As grandezas canônicas adotadas na literatura de engenharia compreendem:

\begin{itemize}
    \item \textbf{Tempo de Subida ($t_r$ --- \textit{Rise Time}):} Intervalo temporal requerido para a resposta transitar de 10\% para 90\% do seu valor final em regime, ou de 0\% a 100\% em sistemas subamortecidos. Mensura a agilidade do controlador.
    \item \textbf{Sobressinal Máximo Percentual ($M_p$ --- \textit{Overshoot}):} Valor de pico máximo atingido pela saída normalizado em relação ao patamar de regime permanente, quantificado percentualmente por:
    \begin{equation}
        M_p = \frac{y_{\max} - y_{\text{final}}}{y_{\text{final}}} \times 100\%
        \label{eq:overshoot}
    \end{equation}
    \item \textbf{Tempo de Acomodação ($t_s$ --- \textit{Settling Time}):} Tempo necessário para que a variável de processo entre e permaneça em uma faixa de tolerância estrita (tipicamente arbitrada em $\pm 2\%$ ou $\pm 5\%$) em torno do valor final.
    \item \textbf{Erro em Regime Permanente ($e_{ss}$ --- \textit{Steady-State Error}):} Diferença residual assintótica entre a referência e a saída quando o tempo tende ao infinito:
    \begin{equation}
        e_{ss} = \lim_{t \to \infty} [r(t) - y(t)]
        \label{eq:erro_regime}
    \end{equation}
\end{itemize}

Para além das métricas pontuais de resposta ao degrau, avalia-se o desempenho dinâmico ao longo de todo o horizonte de ensaio mediante índices integrados de erro e desvio quadrático médio ($\text{RMSE}$):

\begin{align}
    \text{IAE} &= \int_{0}^{T} |e(t)| \, dt \quad \text{(Integral do Erro Absoluto)} \label{eq:iae} \\
    \text{ITAE} &= \int_{0}^{T} t \, |e(t)| \, dt \quad \text{(Integral do Erro Absoluto Ponderado no Tempo)} \label{eq:itae} \\
    \text{RMSE} &= \sqrt{\frac{1}{N} \sum_{k=1}^{N} \left(r[k] - y[k]\right)^2} \quad \text{(Raiz do Erro Quadrático Médio)} \label{eq:rmse}
\end{align}

\subsection{Não-Linearidades Práticas: Saturação e Integrator Windup}
\label{subsec:saturacao_windup}

Em sistemas reais, os atuadores físicos (motores elétricos, servoválvulas, aquecedores) operam sob limites estritos de potência e excursão mecânica. Essa restrição física introduz uma não-linearidade de saturação estática no sinal de controle:

\begin{equation}
    u_{\text{sat}}(t) = 
    \begin{cases}
        u_{\max}, & \text{se } u(t) > u_{\max} \\
        u(t), & \text{se } u_{\min} \le u(t) \le u_{\max} \\
        u_{\min}, & \text{se } u(t) < u_{\min}
    \end{cases}
    \label{eq:saturacao}
\end{equation}

Quando o sinal calculado $u(t)$ excede o limite físico do atuador, o sistema opera transitoriamente em malha aberta no limite de saturação. Nesse intervalo, se o controlador contiver ação integral, o integrador continuará acumulando o erro remanescente ($u_I = K_i \int e \, dt$), crescendo a níveis substanciais. 

Esse fenômeno, denominado \textit{integrator windup}, faz com que, mesmo após o erro inverter de sinal ($y(t) > r(t)$), o atuador permaneça saturado por um período prolongado até que a área negativa do erro compense o valor acumulado no integrador \cite{astrom2006, franklin2015}. As consequências práticas incluem sobressinais desproporcionais, oscilações persistentes e risco de instabilidade física.

Para contornar esse efeito, empregam-se técnicas de compensação \textit{anti-windup}. Dentre as principais abordagens destacam-se o \textit{clamping condicional} (onde a integração numérica é congelada sempre que o atuador saturar com erro de mesmo sinal) e o \textit{back-calculation} (em que a diferença entre o sinal calculado e o sinal saturado é realimentada com ganho $1/T_t$ para descarregar o integrador).

\section{Métodos Clássicos de Sintonia de Controladores}
\label{sec:metodos_sintonia}

A parametrização dos ganhos $K_p, K_i, K_d$ constitui a etapa de sintonia (\textit{tuning}). Em ambientes industriais onde o modelo analítico exato da planta é desconhecido ou complexo demais para ser derivado por equações constitutivas de primeiros princípios, recorre-se a métodos heurísticos e empíricos baseados em ensaios experimentais \cite{cohen1953}.

\subsection{Modelagem por Curva de Reação (Modelo FOPDT)}
\label{subsec:fopdt}

O ensaio da curva de reação consiste em aplicar uma entrada em degrau de amplitude $\Delta u$ à planta operando em malha aberta e registrar a evolução temporal da variável de saída $y(t)$ a partir do repouso. Uma classe expressiva de processos estáveis exibe uma resposta com formato característico em ``S'' (\textit{sigmoide}), que pode ser adequadamente aproximada por um modelo de Primeira Ordem com Tempo Morto (\textit{First Order Plus Dead Time} --- FOPDT):

\begin{equation}
    G(s) = \frac{Y(s)}{U(s)} = \frac{K \, e^{-\theta s}}{\tau s + 1}
    \label{eq:fopdt}
\end{equation}

Os três parâmetros característicos do modelo FOPDT são extraídos pelo método gráfico da reta tangente traçada no ponto de inflexão da curva de resposta (ponto de máxima derivada temporal $\frac{dy}{dt}$), conforme detalhado na Figura \ref{fig:curva_reacao}:

\begin{figure}[htbp]
    \centering
    \caption{Identificação dos parâmetros $K, \tau, \theta$ do modelo FOPDT pelo método da reta tangente no ponto de inflexão da curva de reação.}
    \label{fig:curva_reacao}
    \fbox{\parbox{0.85\textwidth}{\centering
        \vspace{0.4cm}
        \textbf{[Curva de Reação em Degrau $y(t)$]} \\
        \vspace{0.2cm}
        $\Delta y \xrightarrow{} \text{Variação total da saída em regime}$ \\
        $K = \frac{\Delta y}{\Delta u} \quad \text{(Ganho Estático do Processo)}$ \\
        $\theta \quad \text{(Tempo Morto ou Atraso Aparente de Transporte)}$ \\
        $\tau \quad \text{(Constante de Tempo Principal do Sistema)}$
        \vspace{0.4cm}
    }}
    \vspace{0.1cm}
    \fonte{Elaborada pelo autor (2026).}
\end{figure}

\begin{enumerate}
    \item \textbf{Ganho Estático ($K$):} Razão entre a variação total da saída em regime permanente e a amplitude do degrau de entrada aplicado:
    \begin{equation}
        K = \frac{\Delta y}{\Delta u} = \frac{y(\infty) - y(0)}{u(\infty) - u(0)}
        \label{eq:ganho_estatico}
    \end{equation}
    \item \textbf{Tempo Morto Aparente ($\theta$):} Intervalo entre a aplicação do degrau de entrada e a interseção da reta tangente com o valor inicial da saída ($y = 0$).
    \item \textbf{Constante de Tempo ($\tau$):} Intervalo entre o término do tempo morto aparente $\theta$ e o instante em que a reta tangente atinge o patamar de regime permanente $y(\infty)$.
\end{enumerate}

\subsection{Critérios de Sintonia de Ziegler-Nichols em Malha Aberta}
\label{subsec:zn_malha_aberta}

A partir dos parâmetros $K, \tau$ e $\theta$ estimados pelo modelo FOPDT, Ziegler e Nichols \cite{ziegler1942} propuseram fórmulas empíricas de sintonia projetadas para produzir respostas com razão de decaimento de um quarto (\textit{quarter decay ratio}, decaimento na proporção $4:1$ na amplitude das oscilações sucessivas). A Tabela \ref{tab:ziegler_nichols} sintetiza as expressões analíticas para os controladores P, PI e PID na forma padrão e paralela.

\begin{table}[htbp]
    \centering
    \caption{Parâmetros de sintonia clássica de Ziegler-Nichols baseados na curva de reação em malha aberta (modelo FOPDT).}
    \label{tab:ziegler_nichols}
    \begin{tabular}{l|c|c|c|c|c}
        \hline
        \textbf{Tipo} & $K_p$ & $T_i$ & $T_d$ & $K_i = \frac{K_p}{T_i}$ & $K_d = K_p T_d$ \\
        \hline
        \textbf{P}   & $\frac{\tau}{K \theta}$ & --- & --- & --- & --- \\
        \textbf{PI}  & $0{,}9 \frac{\tau}{K \theta}$ & $3{,}33 \, \theta$ & --- & $\frac{K_p}{3{,}33 \, \theta}$ & --- \\
        \textbf{PID} & $1{,}2 \frac{\tau}{K \theta}$ & $2{,}0 \, \theta$ & $0{,}5 \, \theta$ & $\frac{K_p}{2{,}0 \, \theta}$ & $K_p (0{,}5 \, \theta)$ \\
        \hline
    \end{tabular}
    \vspace{0.1cm}
    \fonte{Adaptada de Ziegler e Nichols \cite{ziegler1942} e Ogata \cite{ogata2010}.}
\end{table}

A despeito de sua ampla utilização na indústria, as regras originais de Ziegler-Nichols tendem a gerar respostas relativamente agressivas, com sobressinais da ordem de $20\%$ a $30\%$. Por conseguinte, métodos complementares como o de Cohen-Coon \cite{cohen1953} ajustam essas fórmulas para plantas onde a razão de atraso $\frac{\theta}{\tau}$ é relevante, garantindo maior flexibilidade de parametrização.

\section{Ambientes Virtuais e Simulação no Ensino de Engenharia}
\label{sec:laboratorios_virtuais}

O papel da experimentação prática na formação de engenheiros é amplamente reconhecido na literatura pedagógica \cite{dormido2002, gomes2019}. No contexto da educação em controle e automação, a taxonomia de laboratórios divide-se fundamentalmente em três categorias:

\begin{enumerate}
    \item \textbf{Laboratórios Tradicionais (Físicos):} Baseados em bancadas experimentais físicas locais. Oferecem contato tátil real com sensores, fiação e atuadores, porém impõem custos elevados de aquisição e manutenção, riscos de danos por parametrização incorreta de alunos e restrição de acesso a horários de funcionamento presencial da universidade.
    \item \textbf{Laboratórios Remotos:} Consistem em bancadas físicas reais acopladas a servidores e câmeras de monitoramento, operadas à distância via rede. Preservam a fidelidade física dos experimentos, mas sofrem com limitações de concorrência (um único usuário por bancada física por vez) e dependência de infraestrutura de hardware dedicada.
    \item \textbf{Laboratórios Virtuais (Simuladores):} Ambientes puramente baseados em software nos quais a física do sistema dinâmico é sintetizada por integração numérica em tempo real. Permitem escalabilidade infinita de usuários concorrentes, experimentação irrestrita em ambiente seguro sem risco de colisão física e custo marginal nulo de distribuição.
\end{enumerate}

Para maximizar o engajamento discente, a literatura contemporânea ressalta a importância da \textit{interatividade imersiva} e do \textit{feedback síncrono} \cite{dormido2002}. Interfaces que combinam a escrita de código algorítmico com a visualização tridimensional animada do sistema respondendo em tempo real facilitam a formação de modelos mentais coerentes sobre estabilidade, amortecimento e inércia física.

\section{Tecnologias Computacionais de Base}
\label{sec:tecnologias_computacionais}

A viabilização do \textit{PID Simulator} como laboratório virtual de alta performance executado diretamente no navegador do discente apoia-se em um conjunto de tecnologias modernas de engenharia de software e computação gráfica.

\subsection{Computação Gráfica na Web: WebGL e Three.js}
\label{subsec:threejs}

O padrão WebGL (\textit{Web Graphics Library}) provê uma API JavaScript de baixo nível que permite renderização gráfica tridimensional acelerada por hardware (GPU) diretamente em navegadores web, em conformidade com o padrão OpenGL ES, dispensando a instalação de \textit{plugins} externos \cite{dirksen2018}.

Para abstrair a complexidade de programação de \textit{shaders} em GLSL, matrizes de projeção e gerenciamento de \textit{buffers} de vértices, adotou-se a biblioteca Three.js. Ela estrutura a cena tridimensional em uma hierarquia de nós (\textit{Scene Graph}), gerenciando geometrias poligonais (`BufferGeometry`), materiais físicos com iluminação dinâmica (`SpotLight`, `DirectionalLight`) e câmeras com interpolação de movimento suave (\textit{lerp} --- \textit{Linear Interpolation}) para acompanhamento contínuo dos veículos no circuito.

\subsection{Ambiente Seguro de Execução em Sandbox (Client-Side)}
\label{subsec:sandbox_js}

Permitir que o usuário insira e execute código JavaScript arbitrário em tempo real dentro do navegador acarreta dois desafios fundamentais: estabilidade do laço de renderização e segurança de escopo.

Para mitigar esses riscos sem sobrecarregar servidores remotos com a execução de milhares de rotinas por segundo, implementou-se uma \textit{sandbox} baseada na instanciação isolada de escopo léxico via construtor `new Function(...)`. O código submetido pelo estudante é executado recebendo estritamente os parâmetros do sistema físico ($e[k]$, $\int e$, $\frac{de}{dt}$, $v[k]$, $y[k]$ e $d[k]$), encapsulado em blocos `try/catch` que interceptam erros de sintaxe ou tempo de execução antes que estes alcancem o ciclo de animação `requestAnimationFrame`. Se um erro for identificado, o laço de física é pausado deterministicamente e o painel de telemetria exibe o diagnóstico imediato ao usuário.

\subsection{Arquitetura de Serviços Backend com FastAPI e Docker}
\label{subsec:backend_fastapi}

Para complementar o processamento no cliente, a infraestrutura de retaguarda foi projetada segundo o padrão arquitetural REST \cite{fielding2000}, empregando o framework Python FastAPI. As razões para esta escolha incluem:

\begin{itemize}
    \item \textbf{Alta Performance Assíncrona:} Baseado em ASGI (\textit{Asynchronous Server Gateway Interface}) e `uvicorn`, assegurando tempo de resposta sub-milissegundo para persistência de séries temporais de telemetria;
    \item \textbf{Validação Estrita de Schemas:} Uso do Pydantic v2 para validação e serialização de dados de ensaios laboratoriais e métricas dinâmicas;
    \item \textbf{Processamento Científico e Geração de Relatórios:} Integração nativa com bibliotecas Python (ReportLab, Matplotlib e SciPy) para geração de relatórios técnicos em formato PDF e exportação de datasets científicos em CSV;
    \item \textbf{Conteinerização e Portabilidade (Docker):} Empacotamento integral da aplicação em contêineres Docker padronizados via `docker-compose`, assegurando reproducibilidade dos ensaios em qualquer ambiente operacional.
\end{itemize}

\section{Trabalhos Correlatos e Análise Comparativa}
\label{sec:trabalhos_correlatos}

A literatura científica e o ecossistema educacional apresentam diversas ferramentas destinadas ao ensino de controle. A Tabela \ref{tab:comparativo_ferramentas} estabelece uma comparação sistemática entre as abordagens consolidadas e a plataforma proposta neste trabalho, avaliando seis requisitos essenciais: gratuidade e código aberto, ausência de instalação local, renderização tridimensional interativa, editor de código algorítmico, decomposição didática dos sinais de esforço de controle ($u_P, u_I, u_D$) e módulo de sintonia clássica assistida.

\begin{table}[htbp]
    \centering
    \caption{Quadro comparativo entre o PID Simulator e ferramentas existentes na literatura e no mercado para o ensino de controle.}
    \label{tab:comparativo_ferramentas}
    \resizebox{\textwidth}{!}{%
    \begin{tabular}{l|c|c|c|c|c|c}
        \hline
        \textbf{Ferramenta / Solução} & \textbf{Código Aberto} & \textbf{Web (Sem Instalação)} & \textbf{Cena Física 3D} & \textbf{Editor de Código} & \textbf{Decomposição $u_P, u_I, u_D$} & \textbf{Sintonia Assistida} \\
        \hline
        MATLAB / Simulink & Não & Parcial (Cloud paga) & Parcial (Simscape 3D) & Sim (Scripts) & Parcial & Sim (PID Tuner) \\
        Scilab / Xcos & Sim & Não & Não & Sim & Não & Parcial \\
        Simuladores Web 2D Típicos & Parcial & Sim & Não & Não (Sliders apenas) & Raro & Não \\
        Bancadas Didáticas Físicas & Não & Não & Real (Físico) & Depende do CLP/MCU & Indireta & Manual \\
        \textbf{PID Simulator (Este Trabalho)} & \textbf{Sim} & \textbf{Sim} & \textbf{Sim (Three.js)} & \textbf{Sim (CodeMirror)} & \textbf{Sim (Osciloscópio)} & \textbf{Sim (Z-N / C-C)} \\
        \hline
    \end{tabular}%
    }
    \vspace{0.1cm}
    \fonte{Elaborada pelo autor (2026).}
\end{table}

A análise comparativa evidencia que, embora o ecossistema MATLAB/Simulink ofereça ferramentas avançadas, o custo de licenciamento e a complexidade representam barreiras severas para a massificação discente. Por outro lado, simuladores web existentes restringem-se à manipulação superficial de controles deslizantes (\textit{sliders}) sobre gráficos 2D. 

O diferencial do \textit{PID Simulator} reside na convergência entre acessibilidade universal no navegador, visualização física realista com perturbações dinâmicas de relevo em 3D, escrita de código real pelo estudante e transparência didática completa na instrumentação gráfica multicanal.
